\documentclass[aip,jcp,preprint]{revtex4-1}

\usepackage[utf8]{inputenc}
\usepackage{amsmath}
\usepackage{amssymb}
\usepackage{array}
\usepackage{braket}
\usepackage{booktabs}
\usepackage[capitalize]{cleveref}

\newcommand{\dd}{\mathrm{d}}
\newcommand{\Tr}{\operatorname{Tr}}
\newcommand{\RePart}{\operatorname{Re}}
\newcommand{\ii}{\mathrm{i}}
\newcommand{\xcoord}{{[x]}}
\newcommand{\sfX}{\mathrm{SF\mbox{-}X2C}}
\newcommand{\soX}{\mathrm{SO\mbox{-}X2C}}

\begin{document}

\title{Technical Notes: GHF and X2C-1e Analytic Nuclear Gradients in Forte2}
\author{Forte2 Developers}
\noaffiliation
\date{\today}
\maketitle

\section{Purpose and implementation status}

This note describes the equations and array conventions used to implement
density-fitted generalized Hartree--Fock (GHF) nuclear gradients and
the density-contracted one-electron exact-two-component (X2C-1e) gradient
contribution in Forte2.
The X2C response theory follows Cheng and Gauss,\cite{ChengGauss2011} adapted
to Forte2's decontracted basis, canonical overlap truncation, Pauli-component
ordering, and density-fitting implementation.

The implementation covers the following combinations (SS denotes a
state-specific multiconfigurational wave function):
\begin{center}
\small
\begin{tabular}{@{}>{\raggedright\arraybackslash}p{0.27\linewidth}
                    >{\raggedright\arraybackslash}p{0.49\linewidth}
                    >{\raggedright\arraybackslash}p{0.16\linewidth}@{}}
\toprule
Wave function & One-electron Hamiltonian & Nuclear model \\
\midrule
RHF, UHF, GHF & nonrelativistic & point, Gaussian \\
RHF, UHF, GHF & spin-free X2C-1e & point, Gaussian \\
GHF & spin-orbit X2C-1e & point, Gaussian \\
GHF & spin-orbit X2C-1e with SNSO & point, Gaussian \\
SS-CASSCF/GASSCF & nonrelativistic or spin-free X2C-1e & point, Gaussian \\
2c SS-CASSCF/GASSCF & spinor-upcast nonrelativistic, SF-X2C-1e, or SO-X2C-1e & point, Gaussian \\
\bottomrule
\end{tabular}
\end{center}
All SNSO choices currently recognized by Forte2 (Boettger, DC, DCB, and
row-dependent) are differentiated.  The X2C derivative also follows a fixed
retained overlap eigenspace when canonical orthogonalization removes linear
dependencies.  The number of retained functions must remain unchanged in a
neighborhood of the reference geometry, as is required for any differentiable
truncated-space calculation.

The term ``analytic gradient'' below means that the wave-function derivative
is evaluated from the stationary Lagrangian and the X2C response is propagated
analytically in reverse mode.  With libcint, which is enabled in the default
build, the same-basis $W$ derivative contractions used by X2C are analytic.
Without libcint, \texttt{opVop\_deriv} falls back to a fourth-order finite
difference, so that configuration is not a fully analytic X2C gradient.

\section{Indices, density orientation, and derivative convention}

Greek indices $\mu,\nu,\rho,\sigma$ denote real spatial AO functions, $P,Q$
denote auxiliary functions, $i,j$ denote occupied spinors, and
$\omega,\tau\in\{\alpha,\beta\}$ denote spin components.  A superscript
$\xcoord$ denotes differentiation with respect to one atom-major Cartesian
coordinate $x=R_{Aq}$:
\begin{equation}
 A^\xcoord = \frac{\partial A}{\partial R_{Aq}}.
\end{equation}
Forte2 stores occupied GHF spinors as
\begin{equation}
 C_{\mathrm{occ}} =
 \begin{pmatrix} C^\alpha_{\mathrm{occ}} \\ C^\beta_{\mathrm{occ}}\end{pmatrix},
 \qquad
 C_{\mathrm{occ}}^\dagger
 \begin{pmatrix}S&0\\0&S\end{pmatrix}
 C_{\mathrm{occ}}=1.
\end{equation}
Its density arrays use the orientation
\begin{equation}
 D^{\omega\tau}_{\mu\nu}
 = \sum_i C^\omega_{\mu i} C^{\tau *}_{\nu i}.
 \label{eq:ghf_density_blocks}
\end{equation}
Consequently an AO operator is contracted as
$D_{\nu\mu}h_{\mu\nu}=\Tr(Dh)$.  This orientation matters for complex GHF:
moving the conjugation to the other AO index produces the right real-orbital
limit but the wrong complex gradient.

The spin-traced spatial density is
\begin{equation}
 D_{\mu\nu}=D^{\alpha\alpha}_{\mu\nu}
             +D^{\beta\beta}_{\mu\nu}.
 \label{eq:charge_density}
\end{equation}
It is the only density entering a scalar one-electron operator and the Coulomb
energy.  Off-diagonal spin blocks enter exchange even though they do not enter
\cref{eq:charge_density}.

\section{Density-fitted complex GHF gradient}

\subsection{Metric-inverted three-center tensor}

Let
\begin{align}
 J^P_{\mu\nu} &= (P|\mu\nu), & M_{PQ}&=(P|Q),\\
 Z^P_{\mu\nu} &= (M^{-1})_{PQ}J^Q_{\mu\nu}.
\end{align}
The fitted charge and occupied spinor transition density are
\begin{align}
 \rho^P
 &=D_{\nu\mu}Z^P_{\mu\nu},
 \label{eq:rho}\\
 Q^P_{ij}
 &=\sum_{\omega}\sum_{\mu\nu}
 C^{\omega *}_{\mu i}Z^P_{\mu\nu}C^\omega_{\nu j}.
 \label{eq:q_transition}
\end{align}
The spin sum is inside $Q^P_{ij}$.  This is the GHF distinction from UHF,
where alpha and beta occupied orbitals belong to orthogonal spin sectors and
are treated as two separate occupied lists.

\subsection{Derivative weights}

Differentiating the Coulomb metric and the three-center integrals gives the
weights used by Forte2:
\begin{align}
 W^M_{PQ}
 &=-\frac{1}{2}\rho^P\rho^Q
   +\frac{1}{2}\sum_{ij}Q^P_{ij}Q^Q_{ji},
 \label{eq:ghf_w2}\\
 W^{J,P}_{\mu\nu}
 &=D_{\nu\mu}\rho^P
   -\sum_{\omega ij}
    C^{\omega *}_{\mu i}C^\omega_{\nu j}Q^P_{ji}.
 \label{eq:ghf_w3}
\end{align}
The explicit two-electron derivative is therefore
\begin{equation}
 E^{\xcoord}_{2e}
 =\RePart\left[
 W^{J,P}_{\mu\nu}J^{P,\xcoord}_{\mu\nu}
 +W^M_{PQ}M^{\xcoord}_{PQ}
 \right].
 \label{eq:ghf_df_gradient}
\end{equation}
The derivative integrals are real in the real spherical AO basis, but the
weights may be complex.  Forte2 carries out the complex algebra through
\cref{eq:ghf_w2,eq:ghf_w3} and takes the real part only in the final derivative
contraction.

For canonical occupied spinors, the energy-weighted spinor density is
\begin{equation}
 W_{\mu\nu}^{\omega\tau}
 =\sum_i C^\omega_{\mu i}\epsilon_i C^{\tau *}_{\nu i}.
\end{equation}
Because the AO overlap is spin diagonal, its gradient uses the spin trace
$W^{\alpha\alpha}+W^{\beta\beta}$.  The complete nonrelativistic GHF gradient
is
\begin{align}
 E^\xcoord
 &=E_\mathrm{NN}^\xcoord
 +\RePart\left[D_{\nu\mu}h^\xcoord_{\mu\nu}
 -W_{\nu\mu}S^\xcoord_{\mu\nu}\right]
 +E^{\xcoord}_{2e}.
 \label{eq:ghf_total_gradient}
\end{align}

\section{X2C-1e Hamiltonian used by Forte2}

\subsection{Primitive matrices and Pauli ordering}

X2C is built in a decontracted spatial basis.  The four primitive matrices are
\begin{align}
 S_{\mu\nu}&=\braket{\chi_\mu|\chi_\nu}, &
 T_{\mu\nu}&=\braket{\chi_\mu|\hat p^2/2|\chi_\nu},\\
 V_{\mu\nu}&=\braket{\chi_\mu|\hat V|\chi_\nu}, &
 W_{\mu\nu}&=\braket{\chi_\mu|(\boldsymbol\sigma\cdot\hat{p})
 \hat V(\boldsymbol\sigma\cdot\hat{p})|\chi_\nu}.
\end{align}
The component order used by the public \texttt{opVop} matrices and by the
weights passed to \texttt{opVop\_deriv} is
\begin{equation}
 [W_0,W_x,W_y,W_z]
 =[\hat p\cdot V\hat p,
   (\hat p\times V\hat p)_x,
   (\hat p\times V\hat p)_y,
   (\hat p\times V\hat p)_z].
 \label{eq:opvop_order}
\end{equation}
For SO-X2C these real spatial components are assembled into
\begin{equation}
 W_\mathrm{spinor}=
 \begin{pmatrix}
 W_0+\ii W_z & \ii W_x+W_y\\
 \ii W_x-W_y & W_0-\ii W_z
 \end{pmatrix}.
 \label{eq:pauli_w}
\end{equation}
For SF-X2C only $W_0$ is retained.

\subsection{Orthogonalized Dirac problem}

Let $O$ be the possibly rectangular canonical orthogonalizer of the
decontracted overlap,
\begin{equation}
 O^\dagger S O=1,
\end{equation}
and transform $T,V,W$ to this orthonormal basis.  In a restricted kinetically
balanced basis the generalized one-electron Dirac problem is
\begin{equation}
 \mathcal D C=\mathcal M C\epsilon,
 \label{eq:dirac_problem}
\end{equation}
with
\begin{align}
 \mathcal D&=
 \begin{pmatrix}
 V&T\\T&\dfrac{1}{4c^2}W-T
 \end{pmatrix}, &
 \mathcal M&=
 \begin{pmatrix}
 1&0\\0&\dfrac{1}{2c^2}T
 \end{pmatrix}.
 \label{eq:dirac_matrices}
\end{align}
Partition the positive-energy eigenvectors into large and small rows,
$C_+=(C_L,C_S)^T$.  The decoupling matrix and renormalization metric are
\begin{align}
 X&=C_SC_L^+,\\
 A&=1+\frac{1}{2c^2}X^\dagger T X,
 &R&=A^{-1/2},
 \label{eq:x_and_r}
\end{align}
where $+$ denotes the Moore--Penrose inverse.  The NESC matrix and electronic
Foldy--Wouthuysen Hamiltonian are
\begin{align}
 L&=TX+X^\dagger T-X^\dagger TX+V
   +\frac{1}{4c^2}X^\dagger W X,
 \label{eq:nesc}\\
 h_\mathrm{FW}&=R^\dagger L R.
 \label{eq:hfw}
\end{align}
These are the orthonormal-basis forms of Eqs.~(14)--(17) of
Ref.~\onlinecite{ChengGauss2011}.

\subsection{Back transformation and recontraction}

Forte2 maps \cref{eq:hfw} to the decontracted AO basis with the dual
orthogonalizer $O^{-1}=O^\dagger S$.  It then projects to the original
contracted basis with
\begin{equation}
 P=S_{dd}^{-1}S_{dc},
 \label{eq:projection}
\end{equation}
where $d$ and $c$ label decontracted and contracted bases.  Thus
\begin{equation}
 h_\mathrm{X2C}=P^\dagger (O^{-1})^\dagger
 h_\mathrm{FW}O^{-1}P.
 \label{eq:contracted_hx2c}
\end{equation}
For SO-X2C, $O$, $O^{-1}$, and $P$ are duplicated on the two spin blocks.

\section{Analytic density-contracted X2C derivative}

\subsection{Primitive derivative contractions}

The differential of the X2C energy can be expressed in terms of
\begin{equation}
 S^\xcoord,\quad T^\xcoord,\quad V^\xcoord,\quad
 [W_0^\xcoord,W_x^\xcoord,W_y^\xcoord,W_z^\xcoord].
 \label{eq:x2c_required_derivatives}
\end{equation}
The implementation never materializes these matrices for every coordinate.
Libint2 contracts derivative integrals for $S$, $T$, and point-nucleus $V$
directly against adjoint weights.  Gaussian-nucleus $V^\xcoord$ and all four
$W^\xcoord$ components use libcint blocks that are generated and contracted
one atom at a time.  The relevant libcint kernels are
\texttt{int1e\_ipnuc}, \texttt{int1e\_iprinv},
\texttt{int1e\_ipspnucsp}, and \texttt{int1e\_ipsprinvsp}.

The \texttt{ip} kernels differentiate one AO center.  For atom $A$, Forte2
adds the explicit derivative of the potential centered on $A$, assigns AO
center derivatives to basis functions on $A$, and obtains the second AO-center
term by Hermitian transpose.  The scalar component is symmetric while the
three cross-product components are antisymmetric:
\begin{align}
 (W_0^\xcoord)^T&=W_0^\xcoord, &
 (W_k^\xcoord)^T&=-W_k^\xcoord\quad(k=x,y,z).
\end{align}
The libcint component and transpose conventions are normalized once in the
block generator used by \texttt{opVop\_deriv}; the X2C adjoint supplies weights
only in the public order in \cref{eq:opvop_order}.

\subsection{Orthogonalizer response}

The coordinate-forward differentials in this and the next three subsections
define the linearization whose transpose is applied by the implemented
adjoint.  They are derivational equations; Forte2 does not evaluate or store a
coordinate-forward X2C Hamiltonian derivative.

Inside the retained overlap space the forward differential may be written in
a parallel-transport gauge,
\begin{equation}
 O^\xcoord_\parallel
 =-\frac{1}{2}O(O^\dagger S^\xcoord O).
 \label{eq:orth_deriv_parallel}
\end{equation}
This exactly enforces $(O^\dagger S O)^\xcoord=0$ without singular derivatives
inside degenerate retained eigenspaces.  If overlap eigenvectors $U_d$ have
been discarded and $U_k$ retained, the changing retained subspace adds
\begin{equation}
 O^\xcoord_\perp
 =\sum_{dk}U_d
 \frac{(U_d^\dagger S^\xcoord U_k)}{s_k-s_d}
 s_k^{-1/2}.
 \label{eq:orth_deriv_discarded}
\end{equation}
The derivative of the dual is evaluated from its defining expression,
\begin{equation}
 (O^{-1})^\xcoord=(O^\xcoord)^\dagger S+O^\dagger S^\xcoord.
 \label{eq:dual_deriv}
\end{equation}
For any primitive matrix $B$, its orthonormal-basis derivative is
\begin{equation}
 \bar B^\xcoord=(O^\xcoord)^\dagger BO
 +O^\dagger B^\xcoord O+O^\dagger BO^\xcoord.
 \label{eq:transformed_integral_deriv}
\end{equation}
The method \texttt{\_canonical\_orth\_adjoint} applies the transpose of
\cref{eq:orth_deriv_parallel,eq:orth_deriv_discarded} directly to the overlap
weight.

\subsection{Decoupling response}

From \cref{eq:dirac_matrices}, build $\mathcal D^\xcoord$ and
$\mathcal M^\xcoord$ using $\bar T^\xcoord$, $\bar V^\xcoord$, and
$\bar W^\xcoord$.  Only negative--positive eigenvector coupling is needed for
$X^\xcoord$.  Define
\begin{equation}
 U_{-+}^\xcoord=
 \frac{C_-^\dagger
 (\mathcal D^\xcoord-\mathcal M^\xcoord\epsilon_+)C_+}
 {\epsilon_+-\epsilon_-},
 \label{eq:dirac_response}
\end{equation}
where the division is elementwise over negative and positive states.  Choose
\begin{equation}
 C_+^\xcoord=C_-U_{-+}^\xcoord.
\end{equation}
Positive--positive rotations and normalization changes multiply $C_L$ and
$C_S$ by the same right-hand matrix and therefore cancel from $X=C_SC_L^+$.
With $C_+^\xcoord=(C_L^\xcoord,C_S^\xcoord)^T$,
\begin{equation}
 X^\xcoord=(C_S^\xcoord-XC_L^\xcoord)C_L^+.
 \label{eq:x_deriv}
\end{equation}
Equation \eqref{eq:dirac_response} is the generalized-eigenproblem form of
Eqs.~(31)--(33) of Ref.~\onlinecite{ChengGauss2011}.  It is an uncoupled
one-electron response, not an iterative CPHF equation.  In the implementation,
the negative--positive denominators are applied in reverse to accumulate
adjoints of $\mathcal D$ and $\mathcal M$.

\subsection{Renormalization and Foldy--Wouthuysen response}

Differentiate the metric in \cref{eq:x_and_r}:
\begin{equation}
 A^\xcoord=\frac{1}{2c^2}\left[
 (X^\xcoord)^\dagger TX+X^\dagger T^\xcoord X
 +X^\dagger T X^\xcoord\right].
\end{equation}
For $A=U\,\mathrm{diag}(a)U^\dagger$, the inverse-square-root Frechet
derivative defining the linearization is
\begin{align}
 R^\xcoord
 &=U\left[F\circ(U^\dagger A^\xcoord U)\right]U^\dagger,\\
 F_{ij}
 &=\begin{cases}
 \dfrac{a_i^{-1/2}-a_j^{-1/2}}{a_i-a_j},&a_i\ne a_j,\\[6pt]
 -\dfrac{1}{2}a_i^{-3/2},&a_i=a_j.
 \end{cases}
 \label{eq:r_frechet}
\end{align}
The equal-eigenvalue branch makes the expression stable for degeneracies.
This Frechet map is self-adjoint on Hermitian matrices, so the reverse response
uses the same divided-difference kernel.

Direct differentiation of \cref{eq:nesc} gives
\begin{align}
 L^\xcoord={}&T^\xcoord X+TX^\xcoord
 +(X^\xcoord)^\dagger T+X^\dagger T^\xcoord
 -(X^\xcoord)^\dagger TX-X^\dagger T^\xcoord X-X^\dagger TX^\xcoord
 +V^\xcoord \nonumber\\
 &+\frac{1}{4c^2}\left[(X^\xcoord)^\dagger WX
 +X^\dagger W^\xcoord X+X^\dagger WX^\xcoord\right],
 \label{eq:l_deriv}\\
 h_\mathrm{FW}^\xcoord={}&
 (R^\xcoord)^\dagger LR+R^\dagger L^\xcoord R+R^\dagger LR^\xcoord.
 \label{eq:hfw_deriv}
\end{align}
Omitting $X^\xcoord$ or $R^\xcoord$ drops picture-change response and does not
give the derivative of the many-electron X2C-1e energy.\cite{ChengGauss2011}

\subsection{Projection and SNSO derivatives}

The decontracted-to-contracted projection response follows from
\cref{eq:projection}:
\begin{equation}
 P^\xcoord=S_{dd}^{-1}
 \left(S_{dc}^\xcoord-S_{dd}^\xcoord P\right).
 \label{eq:projection_deriv}
\end{equation}
Applying \cref{eq:dual_deriv,eq:projection_deriv} to
\cref{eq:contracted_hx2c} defines the corresponding forward differential.  The
code instead propagates the contracted density backward through these maps to
obtain weights for $S_{dd}^\xcoord$ and $S_{dc}^\xcoord$.

SNSO factors in Forte2 depend on atomic number and shell angular momentum, not
on geometry.  The code decomposes the primitive SO Hamiltonian into Pauli
components and applies the same self-adjoint elementwise SNSO map to the primal
Hamiltonian and its adjoint.  There is therefore no response of the SNSO factors
themselves.

\subsection{Density-contracted adjoint used in production}

Forte2 exposes only the density-contracted X2C gradient contribution; there is
no matrix-valued $h_\mathrm{X2C}^\xcoord$ API.  The code propagates the AO
density backward through the projection, SNSO map, dual orthogonalizer,
Foldy--Wouthuysen transformation, inverse square root, NESC matrix, decoupling
response, Dirac eigensystem, and canonical orthogonalizer.
It produces five primitive adjoints,
\begin{equation}
 \overline S,\quad \overline T,\quad \overline V,\quad
 [\overline W_0,\overline W_x,\overline W_y,\overline W_z],\quad
 \overline S_{dc},
\end{equation}
such that
\begin{equation}
 E_{1e,\mathrm{X2C}}^\xcoord
 =\RePart\left[
 \overline S:S^\xcoord+\overline T:T^\xcoord
 +\overline V:V^\xcoord+\sum_k\overline W_k:W_k^\xcoord
 +\overline S_{dc}:S_{dc}^\xcoord\right].
 \label{eq:x2c_adjoint_contraction}
\end{equation}
Here $A:B$ denotes the elementwise contraction in Forte2's stored matrix
orientation.  The inverse-square-root Frechet map in
\cref{eq:r_frechet} propagates the renormalization adjoint with the same
divided-difference kernel.  The
negative--positive energy denominators in \cref{eq:dirac_response} likewise
appear once in the reverse decoupling response.

Libint2 contracts primitive adjoints directly in its compiled derivative
kernels.  For Gaussian-nucleus $V^\xcoord$ and analytic libcint
$W^\xcoord$, one-atom derivative blocks are returned and immediately contracted
in Python.  Consequently the production path does not allocate any array of
shape $(3N_\mathrm{atom},n,n)$ or
$(4,3N_\mathrm{atom},n,n)$.  Its dense X2C matrix algebra is
$O(n^3)$ rather than repeating $O(n^3)$ response work for every nuclear
coordinate; the unavoidable derivative-integral contractions remain.

\section{Variational gradient assembly with X2C}

X2C-1e transforms only the one-electron Hamiltonian.  The AO overlap and
two-electron integrals remain untransformed.  Accordingly the X2C SCF gradient
is obtained from \cref{eq:ghf_total_gradient} by replacing the nonrelativistic
$T^\xcoord+V^\xcoord$ contraction with
\begin{equation}
 E_{1e,\mathrm{X2C}}^\xcoord
 =\RePart\Tr(D_\mathrm{hcore}h_\mathrm{X2C}^\xcoord).
 \label{eq:x2c_hcore_contraction}
\end{equation}
For RHF, UHF, and SF-X2C GHF, $D_\mathrm{hcore}$ is the spatial spin trace.
For SO-X2C GHF it is the full complex spinor density.  The overlap Pulay term
still uses the spatial spin trace of the energy-weighted density, and the DF
two-electron terms still use \cref{eq:ghf_w2,eq:ghf_w3}.

The same replacement applies to a fully variational, single-state CASSCF or
GASSCF energy.  For a two-component calculation, collect the inactive core
and active spinors into $C_h=[C_i\ C_u]$ and define
\begin{equation}
 \gamma^h =
 \begin{pmatrix}I_{\mathrm{core}}&0\\0&\gamma^{\mathrm{act}}\end{pmatrix},
 \qquad
 D_{\mathrm{spinor}}=C_h(\gamma^h)^T C_h^\dagger.
 \label{eq:rel_casscf_density}
\end{equation}
The transpose in \cref{eq:rel_casscf_density} is required by the Forte2
relativistic RDM convention: MO integrals and RDMs contract element by element,
rather than through an implicit matrix transpose.  The spatial density used by
scalar operators is $D^{\alpha\alpha}+D^{\beta\beta}$.

The active spin-orbital cumulant is
\begin{equation}
 \lambda_{uv,wx}
 =\Gamma_{uv,wx}-\gamma_{uw}\gamma_{vx}
  +\gamma_{ux}\gamma_{vw}.
\end{equation}
With $Z^P_{xy}=\sum_\omega
(C_h^\omega)^\dagger Z^P C_h^\omega$, the compact CASSCF DF weights have the
same form as the spin-free weights, except that the exchange coefficient is
one instead of one half.  Complex three-center adjoints are returned to the
spatial AO basis as
\begin{equation}
 W^P_{\mu\nu}=\sum_{xy\omega}
 C^{\omega *}_{\mu x}W^P_{xy}C^\omega_{\nu y}.
 \label{eq:rel_casscf_w3_backtransform}
\end{equation}
Using $C W C^\dagger$ in
\cref{eq:rel_casscf_w3_backtransform} is valid only for real orbitals and
breaks arbitrary spinor-phase invariance.  The Pulay term is obtained from the
Hermitian part of the converged relativistic orbital Lagrangian and then spin
traced in the AO basis.  No CASSCF Z-vector is needed for the optimized
single-state objective when every nonredundant orbital rotation is variational.

This implementation change concerns only the X2C one-electron term.  The
density-fitting tensors $Z^P_{\mu\nu}$, $W^M_{PQ}$, and
$W^{J,P}_{\mu\nu}$ remain materialized by the existing SCF and
multiconfigurational gradient code.

For all supported SCF and multiconfigurational gradients, Gaussian nuclei
change three pieces consistently: $V^\xcoord$, $W^\xcoord$, and
$E_\mathrm{NN}^\xcoord$.  The last is
evaluated by contracting the derivative of the Coulomb interaction between
normalized Gaussian nuclear charge functions, with diagonal self-interactions
removed exactly as in the energy.

\section{Code map}

\begin{center}
\begin{tabular}{p{0.35\linewidth}p{0.57\linewidth}}
\toprule
File or function & Responsibility \\
\midrule
\texttt{forte2/scf/ghf\_grad.py} & GHF densities, energy-weighted density, and final driver.\\
\texttt{forte2/scf/hf\_grad.py} & Shared RHF/UHF DF weights and complex GHF DF weights.\\
\texttt{mcopt/mc\_optimizer\_grad.py} & Spin-free and spinor CASSCF densities, DF weights, and final driver.\\
\texttt{mcopt/orbital\_optimizer.py} & Real and two-component orbital Lagrangians for the Pulay term.\\
\texttt{forte2/gradients/utils.py} & Common one-electron, overlap, DF, and nuclear-repulsion assembly.\\
\texttt{one\_electron\_deriv*} & Libint2 contractions of $S^\xcoord$, $T^\xcoord$, and point-$V^\xcoord$ with adjoint weights.\\
\texttt{integrals.py} & Contracted Gaussian-$V^\xcoord$ and full $W^\xcoord$, component normalization, and backend selection.\\
\texttt{libcint\_two\_center.h} & Bound libcint derivative kernels.\\
\texttt{forte2/x2c/x2c.py} & X2C Hamiltonian and density-contracted gradient entry point.\\
\texttt{forte2/x2c/x2c\_grad.py} & Density-contracted reverse response and direct primitive derivative contractions.\\
\bottomrule
\end{tabular}
\end{center}

\section{Verification strategy}

The implementation is tested in layers so that a failure identifies the
responsible algebra:
\begin{enumerate}
\item Contracted Libint2 $S^\xcoord$, $T^\xcoord$, and point-$V^\xcoord$
contributions are compared directly with finite differences.
\item Contracted Gaussian-nucleus $V^\xcoord$ and all four $W^\xcoord$
components are compared with finite differences using fixed random weights.
\item The X2C adjoint is compared with fourth-order finite differences of
$\RePart\Tr(Dh_\mathrm{X2C})$ for fixed Hermitian real and complex densities.
The cases cover SF-X2C, SO-X2C, both nuclear models, all SNSO variants, and a
discarded overlap eigenvector.
\item Nonrelativistic complex GHF is compared with energy finite differences
and with the collinear UHF limit.
\item Full SF-X2C UHF and SO-X2C GHF gradients are compared with finite
differences of converged SCF energies and checked for translational invariance.
\item Two-component CASSCF/GASSCF gradients are checked against the
nonrelativistic upcast limit with arbitrary spinor phases and against SF- and
SO-X2C energy finite differences.  The tests include a correlated triatomic
SNSO CASSCF calculation and an SNSO GASSCF calculation with optimized inter-GAS
rotations.  Gaussian nuclei are checked for nonrelativistic CASSCF and
SO-X2C/SNSO GASSCF.
\end{enumerate}

Useful diagnostics are agreement of the fixed-density X2C contraction with
finite differences, zero total translation of an isolated molecule, arbitrary
spinor-phase invariance, and agreement of the complex GHF collinear limit with
UHF.  A significant imaginary part before the final real projection diagnoses
loss of Hermiticity or a conjugation error; it should not be discarded early to
conceal such an error.

\section{Known boundaries}

The current gradient implementation assumes density fitting, not Cholesky
two-electron integrals.  X2C $W^\xcoord$ is fully analytic when libcint is
available, which is the default Forte2 build.  Without libcint, point-nucleus
$W^\xcoord$ contractions use the fourth-order fallback and Gaussian nuclear
models are unavailable.  Derivative integrals using Libint2 remain subject to
the angular-momentum limit of the installed Libint2 library.  A geometry at
which the number of retained overlap eigenvectors changes is
non-differentiable under hard canonical truncation and must be handled by
changing the threshold or basis.

CASSCF/GASSCF gradients are currently limited to a single optimized state and
orbital spaces without frozen core, frozen virtual, active-frozen, or frozen
inter-GAS rotations.  Both the orbital optimization and CI roots must be
converged before a gradient is evaluated.  The excluded cases require response
terms outside the stationary single-state Lagrangian used here.

The implemented Hamiltonian is X2C-1e: two-electron picture-change corrections
are not included in either the energy or gradient.  This is the same
approximation as the existing Forte2 X2C energy and is distinct from X2C-2e or
a four-component Dirac--Coulomb gradient.

\begin{thebibliography}{9}
\bibitem{ChengGauss2011}
L. Cheng and J. Gauss,
``Analytic energy gradients for the spin-free exact two-component theory using
an exact block diagonalization for the one-electron Dirac Hamiltonian,''
J. Chem. Phys. \textbf{135}, 084114 (2011).
\end{thebibliography}

\end{document}
